La comprensión de los mecanismos internos mediante los cuales un programa escrito en un lenguaje de alto nivel se transforma en un cómputo ejecutable constituye uno de los pilares formativos de la Ingeniería de Sistemas. Los cursos de interpretación y compilación de lenguajes de programación introducen al estudiante en procesos como el análisis léxico y sintáctico, la construcción de árboles de sintaxis abstracta y la gestión de ambientes de ejecución, conceptos que sustentan áreas tan diversas como el desarrollo de compiladores, las herramientas de análisis estático y el diseño de lenguajes de dominio específico. Sin embargo, la naturaleza abstracta de estos contenidos y su distancia respecto a la experiencia de programación habitual del estudiante hacen de su enseñanza un reto reconocido en la literatura de educación en computación \cite{liu2017evaluation,navas2022modular}.

En el programa de Ingeniería de Sistemas de la Universidad del Valle sede Tuluá, este reto se concentra en la asignatura \textit{Fundamentos de Interpretación y Compilación de Lenguajes de Programación} (FLP), que emplea el lenguaje Racket junto con los materiales del libro \textit{Essentials of Programming Languages} (EOPL) como entorno principal de trabajo. Si bien esta combinación ofrece una base conceptual rigurosa, exige del estudiante una familiaridad previa con el paradigma funcional y con una sintaxis que, en ausencia de dicha familiaridad, tiende a desplazar su atención desde los conceptos de interpretación hacia la resolución de errores sintácticos \cite{Spigariol2023Aprendiendo}. A esta fricción se suma la carencia de recursos interactivos que permitan explorar de forma autónoma los procesos internos estudiados, pese a que la evidencia disponible señala que las representaciones visuales e interactivas favorecen la comprensión de estos procesos al hacer tangibles las transformaciones que ocurren en cada etapa del procesamiento de un programa \cite{jovanovic2021comvis,lu2022usability,Stamenkovic2024AWE}.

Frente a este panorama, el presente trabajo de grado propone el diseño, la implementación y la evaluación de FLP Viewer, un prototipo web orientado a apoyar la comprensión de los contenidos de la asignatura FLP mediante la visualización de las estructuras internas que un intérprete construye y manipula durante la evaluación de un programa. A diferencia de las herramientas de visualización revisadas en el estado del arte, que operan sobre lenguajes predefinidos, el prototipo permite al estudiante definir su propia gramática mediante notación BNF y observar, de forma simultánea, el árbol de sintaxis abstracta generado y la evolución del ambiente de ejecución asociado.

El Capítulo~\ref{cap1} formula el problema de investigación con base en el diagnóstico realizado, plantea los objetivos del trabajo y delimita su alcance. Los capítulos posteriores documentan el marco referencial, el diagnóstico de dificultades conceptuales, el diseño y la implementación del prototipo, y los resultados de su evaluación con estudiantes de la asignatura.
